\documentclass[12pt,a4paper]{article}
\usepackage[utf8]{inputenc}
\usepackage{graphicx}
\usepackage{color}
\usepackage[magyar]{babel}
\usepackage{listings}

\usepackage{enumerate}
\usepackage{tikz}
\usetikzlibrary{shapes,arrows}
\usetikzlibrary{positioning}

\renewcommand{\lstlistingname}{Lista}
\usepackage{hyperref}
\hypersetup{
	pdftitle={HA5KFU tanfolyam - Forgalmazási ismeretek},
	pdfauthor={HA5KFU Rádióamatőr Klub},
	pdfsubject={HA5KFU tanfolyam},
	pdfcreator={latex},
	pdfkeywords={ },
	pdfpagemode=UseOutlines,
	pdfdisplaydoctitle=true,
	pdflang=hu,
	unicode
}

\pagestyle{plain}
\sloppy
\begin{document}
\begin{center}
\includegraphics[width=300pt,keepaspectratio]{figures/ha5kfu.eps}
\\[0.5cm]
Rádióamatőr tanfolyamot segítő jegyzet, egyelőre kidolgozás alatt \\
Összeállította: Szabó Áron % Feel free to add yourself
\\[1cm]

{\huge \bfseries Forgalmazási ismeretek \\[2cm]}



\end{center}

\renewcommand{\contentsname}{Tartalom}\tableofcontents 
\newpage
\section{Jogszabályok}

A rádióamatőr tevékenységet a világ szinte minden országában jogszabályok szabályozzák. Magyarországon a hazai törvények és  \textbf{NMHH} (Nemzeti Média és Hírközlési Hatóság) rendeletek.

Többek között a rádióamatőr vizsgáról, engedélykiadásról és forgalmazásról rendelkeznek ezek a szabályozások.  

A rádióamatőr szolgálatra nemzetközi szabályok is érvényesek. A Nemzetközi Távközlési Egyesület (International Telecommunication Union,  \textbf{ITU}) osztja fel a frekvenciatartományt sávokra, így az amatőr frekvenciasávokat illetve ezeken a szolgálat jogállását (elsődleges vagy másodlagos). Az ITU három régióra osztja fel a Földet, egy régión belül a szabályozások egységesek \cite{iaru_regi}.

A Nemzetközi Rádióamatőr Szövetség (International Amateur Radio Union, \textbf{IARU}) tesz ajánlást az amatőr sávok felosztására. Általában a sáv elején a kisebb (CW, SSB), a végén a nagyobb sávszélességű módok (FM) találhatók. A sávfelosztás ajánlást ad az átjátszóállomások és beacon-ök frekvenciáira is \cite{iaru}. 

Európában és néhány más országban a Postai és Távközlési Igazgatások Európai Értekezlete (European Conference of Postal and Telecommunications Administrations, \textbf{CEPT}) kialakított egy egységesített vizsgarendszert (Harmonized Amateur Radio Examination Certificate, HAREC), hogy a tagállamok ezeket a vizsgákat kölcsönösen elismerjék, az alapján rádióamatőr engedélyt adjanak ki. 
\subsection{Rádióamatőr engedély}
Magyarországon vizsgabizonyítvány birtokában az NMHH állítja ki a rádióamatőr engedélyt. Amatőr rádióállomást üzemeltetni, önállóan forgalmatni csak amatőrengedély birtokában szabad.

Háromféle amatőrállomás lehet:

\begin{itemize}
\item Egyéni: az amatőr engedély egy természetes személyhez tartozik
\item Közösségi: egy rádióamatőr klub tartja üzemben
\item Különleges amatőrállomás: pl. felügyelet nélküli jeladó, vagy átjátszó állomás, alkalmi amatőrállomás
\end{itemize} 

Az egyéni amatőr engedélynek a vizsga szintje alapján háromféle fokozata van: 

\begin{itemize}
\item Kezdő (csak 14 és 18 éves kor között, illetve 60 éves kor felett igénylehtő)
\item CEPT Novice
\item CEPT HAREC
\end{itemize} 
Az alacsonyabb szintű engedélyeken teljesítmény és egyéb korlátozások alá esnek. Közösségi állomás csak CEPT HAREC egyéni engedély birtokában üzemeltethető.

A rádióamatőr engedély tartalmazza a Morze engedélyt is (ha van), távíró módot csak sikeres Morze vizsgával lehet használni.

\subsubsection{Hívójelek}

A hívójel segítségével a világ bármelyik amatőrállomása azonosítható, angol betűkből és számjegyekből áll. Két része az előtag és az utótag.

Az előtag első néhány karaktere az (ITU által kijelölt) országkód, Magyarország esetén HA vagy HG. Az országkód utáni részt a nemzetek saját hatáskörben oszthatják ki, hazánkban ez egy számjegy (ami régebbi szabályozás szerint Magyarország egy régióját azonosította).

Az utótag Magyarországon szabadon választható. A teljes hívójel (egyéni és közösségi) minimum 5, maximum 7 karakterből állhat.

\subsubsection{QTH}

A rádióamatőr engedély tartalmazza az amatőrállomás telepítési helyét is (rövidítése QTH). Ha ideiglenesen a telepítési helyen kívülről forgalmazunk, használhatjuk a hívójelben a /P (portable) jelet a hívójelünk után.

\section{Forgalmazás}

A rádióamatőr másik rádióamatőrrel létesíthet összeköttetést, ezt nevezzük forgalmazásnak. A forgalmazást (amatőrök számára) közérthető nyelven, titkosítatlanul kell folytatni, tartalma szakmai (tilos politika, ipari, gazdasági adatok, megtévesztő jelek, műsor sugárzása).

Az amatőr rádióadást hívójellel kell azonosítani.

A forgalmazást sokféle üzemmódban lehet folytatni, néhány példa:

\begin{itemize}
\item CW: Continuous Wave, távíró üzemmód. A vivőfrekvencia megléte vagy hiánya adja a két állapotú jelet. A Morze ABC segítségével tudunk karaktereket átvinni. A CW-nek van a legkisebb sávszélesség igénye az összes üzemmód közül.
\item Fónia: a rádióamatőr beszéddel kommunikál
\begin{itemize}
\item SSB: Single Side Band AM: Amplitúdómodulált jel, amiből a vivő és az egyik oldalsáv elnyomott. Így kisebb teljesítményen és kisebb sávszélességen lehet adni. Az FM-nél gyengébb hangminőségű, de kevésbé érzékeny a zajra. Rövidhullámon ez az üzemmód jellemző, távolsági összeköttetésekre (DX) a legalkalmasabb.
\item FM: Frekvenciamodulált jel, használata az ultrarövid amatőrsávon jellemző.
\item Digitális: a hang átvitelére valamilyen digitális kódolást használ. Jó hangminőség, de érzékeny a zajhatásokra.
\end{itemize}
\item Karakteres átvitel
\item Képtovábbítás (Slow Scan Television, ATV) 
\end{itemize}
\subsection{Fónia forgalmazás}

Amikor egy frekvencián el szeretnénk kezdeni adni, figyeljünk be, hogy azt nem használja-e már valaki más. Kérdezzünk be: "Szabad ez a frekvencia?", ha nem válaszolnak, biztosra vehetjük, hogy nem zavarunk olyat, akit mi nem hallunk.

Ha bármilyen másik amatőrállomással szeretnénk kapcsolatot létesíteni, adjunk általános hívást (angolul CQ). Ezt és minden további adást is azonosítsunk a hívójelünkkel.

Adás során érdemes a fontosabb üzeneteket megismételni, hogy a másik fél biztosan értse. A nehezen érthető szavakat, hívójelet érdemes betűzni. Az adás végét zárjuk le: "Vissza hozzád" (angolul Back to you), hogy a másik tudja, hogy mikor adhat, illetve ne zavarjon bele az adásunkba (szimplex módban adás közben nincs vétel).

\subsection{Riport}

A rádióamatőrök az adást riporttal értékelik, morze esetén RST, fónia esetén RS.

\begin{itemize}
\item R: \textbf{Readability} (1-5) Érthetőség, mennyire érthetően, tisztán hallottuk a másikat
\item S: \textbf{Signal} (1-9) Jelszint, a rádiónk kijelzőjéről tudjuk leolvasni
\item T: \textbf{Tone} (1-5) A CW jel minősége
\end{itemize}

Egy tisztán érthető, nagy jelszintű fónia adásra például 59-es (ötös kilences) riportot adunk.

\subsection{Naplózás}

Az összeköttetésekről naplót kell vezetni, amiben az alábbi információk szerepelnek:
\begin{itemize}
\item \textbf{Sorszám}
\item \textbf{Dátum és idő}: a forgalmazás kezdetének ideje GMT időzóna szerint
\item \textbf{Sáv vagy frekvencia} ahol az összeköttetés történt
\item \textbf{Az ellenállomás hívójele}
\item \textbf{Adott és vett riport}, pl 59 - 59 ha kifogástalan az összeköttetés
\item \textbf{Üzemmód}: FM, SSB...
\item Egyéb információk: QTH (ellenállomás helye) vagy QRA (az ellenállomás /rádióklub/ neve), az operátor neve, stb.
\end{itemize}

A naplót elsősorban saját célra vezetjük, de rádióamatőr versenyeken is ezzel hitelesítjük az összeköttetést (sorszámmal), és a hatóság is elkérheti.

\subsection{Szakzsargon}

Mint a legtöbb szakterületnek, a rádiózásnak is vannak különféle szakkifejezései, amik közül néhányat összegyűtöttünk:

\begin{itemize}
\item \textbf{QTH:} rádióállomás telepítési helye (pl. "QTH-m Budapest")
\item \textbf{QSO:} közvetlen összeköttetés (pl. "ma QSO-ztam egy japán állomással")
\item \textbf{QSL:} vétel nyugtázása. Rádióamatőröknél kultúrája van QSL-lapok küldésének, amik általában képeslap-szerű kártyák, és pl. a nagytávolságú összeköttetések emlékére küldik egymásnak az állomások.
\item \textbf{DX:} nagytávolságú összekköttetés (általában rövidhullámon, másik országbeli állomással)
\item \textbf{CQ:} minden állomásnak szóló hívás ("seek you") - egy frekvencián CQ-t adunk, és várjuk, hogy egy másik állomás válaszoljon. 
\item \textbf{QRM, QRN:} Mesterséges, vagy természetes légköri zavar jelzése
\item \textbf{CW:} continuous wave (folyamatos hullám) adásmód, morze kódot adhatunk ebben az üzemmódban
\item \textbf{DE:} -tól, -től. HA8KDA DE HA1FLX = HA8KDA-nak HA1FLX-től
\item \textbf{73:} elköszönés (minden jót), ezt a kódot táviratoknál használták, és a rádióamatőr kultúrában megmaradt


\end{itemize}

Nemzetközi összeköttetéseknél hívójelek, Q-kódok stb. betűzésére a NATO betűzési ábécét használjuk, amit minden rádióamatőrnek ismernie kell \cite{nato}.

\subsection{Rádióamatőr sávok}

A rádióamatőrök csak a számukra kijelölt sávokon, más rádióamatőrökkel forgalmazhatnak. Nagyrészt minden frekvenciatartományban van rádióamatőröknek kijelölt sáv, így az amatőrök tudnak kísérletezni az egyes frekvenciasávok tulajdonságaival.


\subsubsection{Amatőr sávok felosztása}

A frekvenciakiosztás az IARU honlapján elérhető \cite{iaru}.

\subsection{Zajzár}

A zajzár feladata, hogy a vett jelet csak adás esetén halljuk: ha csak zaj, vagy számunkra érdektelen adást veszünk, a zajzár bekapcsol, és lenémítja a rádiót.

\paragraph{Squelch}

Squelch (rövidítve SQL), magyarul zajzár, ha egy bizonyos jelerősség alatt van a vétel, a hangszóró kikapcsol. Ez FM üzemmód esetén hasznos, ahol hangos zajt hallunk, ha nem hangolunk adásra.
\paragraph{CTCSS}

A Continuous Tone-Coded Squelch System (CTCSS) egy alacsony frekvenciájú színusz jelet ad a hanghoz, aminek 50 különböző frekvenciája lehet. Egy flottához tartozó rádiók ugyanazt a frekvenciájú CTCSS-t használják, így a zajzár csak akkor nyit ki, ha az adás mellé a beállított frekvenciájú hullám van hozzáadva.

Fontos, hogy a CTCSS nem véd az interferencia ellen, így az egy frekvenciára hangolt, különböző CTCSS-t használó rádiók ugyanúgy zavarhatják egymást (csak egymás adását nem veszik észre).
\paragraph{DCS}
A Digital-Coded Squelch hasonlóan működik a CTCSS-hez, csak az folyamatos hullám helyett digitális bitsorozatot használ.
\subsection{Átjátszók}
Átjátszóállomásokat közvetlen terjedés esetén használunk, amikor nem látunk rá a másik állomásra. Az átjátszónak van egy megadott felmenő és lejövő frekvenciája, ezeken veszi és adja az átjátszott adásokat. Ezek nem lehetnek azonosak, hiszen akkor az átjátszó a vett adását zavarná, és azonnal begerjedne. Az átjátszó használatához rádiónkat úgy kell beállítani, hogy a lejövő frekvencián vegyen, de adáshoz kapcsoljon át a felmenő frekvenciára. 

Egyes átjátszóállomások úgy vannak beállítva, hogy csak a megfelelő CTCSS vagy DCS kód észlelésekor lépnek működésbe.
\newpage

\renewcommand{\refname}{Irodalomjegyzék}
\begin{thebibliography}{9}
\bibitem{5kdr_j} 
Budapest Fővárosi Rádióamatőr Klub 
\textit{Rádióamatőr vizsgára felkészítő jegyzet}. 
\\\url{https://www.ha5kdr.hu/}

\bibitem{rendelet} 
\textit{15/2013. (IX. 25.) NMHH rendelet a rádióamatőr szolgálatról} 
\\\url{http://njt.hu/cgi_bin/njt_doc.cgi?docid=163445}

\bibitem{iaru_regi} 
\textit{IARU / ITU Regions} 
\\\url{http://www.iaru.org/regions.html}

\bibitem{iaru} 
\textit{IARU Spectrum and Band Plans} 
\\\url{https://www.iaru-r1.org/index.php/spectrum-and-band-plans/}


\bibitem{nato} 
\textit{NATO phonetic alphabet} 
\\\url{https://www.nato.int/alphabet}

\end{thebibliography}

\end{document}
